% =====================================================================
% Capitolo 3: integrità e checksum
% =====================================================================
% copre: docs/03-integrita-checksum.md
% verificato-al-commit: 3f1c9b3
% ---------------------------------------------------------------------

\chapter{Come un sistema si accorge che un dato è rotto}
\label{cap:checksum}

\section{Il problema, prima della soluzione}

Un calcolatore che salva dei dati su un supporto e poi li rilegge ha un problema che a
prima vista non sembra esistere: non ha modo di sapere se ciò che rilegge sia ciò che
aveva scritto. Il supporto può guastarsi, l'alimentazione può mancare a metà della
scrittura, una batteria può esaurirsi, e in nessuno di questi casi i dati diventano
visibilmente strani. Diventano soltanto diversi.

Vale la pena fermarsi su questo punto, perché è l'origine di tutto il capitolo. Un
numero sbagliato ha esattamente lo stesso aspetto di un numero giusto. Se un byte che
valeva \hx{2A} adesso vale \hx{2B}, non c'è nulla in quel byte che lo denunci: sono
otto bit, erano otto bit prima e sono otto bit adesso, e ogni combinazione di otto bit
è un valore legittimo. Il dato non porta con sé l'informazione di essere corretto.

La conseguenza pratica, su un gioco che salva su una cartuccia, è che un guasto
silenzioso è peggio di un guasto rumoroso. Un salvataggio che non si carica affatto è
un problema evidente; un salvataggio che si carica con un valore sbagliato in un punto
qualunque è un problema che si scopre trenta ore dopo, quando le conseguenze si sono
propagate. La difesa contro il secondo caso è l'argomento di questo capitolo.

\section{L'idea: scrivere due volte la stessa cosa, in due modi diversi}

Se il dato non può dire di sé stesso di essere corretto, gli si affianca un secondo
dato che lo riassume. Al momento della scrittura si calcola quel riassunto e lo si
memorizza accanto ai dati; al momento della lettura lo si ricalcola sui dati riletti e
lo si confronta con quello memorizzato. Se i due coincidono, i dati sono probabilmente
integri. Se differiscono, qualcosa è cambiato.

Questo riassunto si chiama \term{checksum}, che alla lettera vuol dire somma di
controllo, e nella forma più semplice è proprio una somma.

Prima di guardare come i giochi lo fanno, conviene essere precisi su cosa un checksum
può e non può fare, perché è la parte che si fraintende più spesso.

Un checksum \emph{rileva} un errore, non lo \emph{corregge}. Sapere che la somma non
torna dice che almeno un bit è cambiato, non quale: da un riassunto non si ricostruisce
il testo. Correggere richiede una ridondanza molto maggiore, cioè i codici a correzione
d'errore, che questi giochi non usano perché costerebbero spazio che non avevano.

Un checksum, inoltre, \emph{può sbagliarsi}. Poiché comprime molti byte in pochi, esistono
necessariamente configurazioni diverse dei dati che producono lo stesso riassunto: è
una conseguenza del principio dei cassetti, non un difetto dell'algoritmo. Un checksum a
16 bit ha 65\,536 valori possibili, quindi una corruzione casuale ha circa una
possibilità su 65\,536 di produrre per caso il valore giusto e passare inosservata. Un
checksum a 8 bit ne ha 256, e la probabilità sale a circa una su 256.

Detto in termini che a chi viene dalle telecomunicazioni suonano familiari, un checksum
è un codice di rilevazione d'errore con potere di rilevazione limitato e capacità di
correzione nulla, e la lunghezza del checksum è il parametro che governa la
probabilità di errore non rilevato.

\section{La forma più semplice: la somma degli otto bit}

La generazione 1, cioè Pokemon Rosso, Blu e Giallo, protegge l'area principale del suo
salvataggio con un checksum a un solo byte, e la regola è la più semplice
immaginabile: si sommano tutti i byte dell'area protetta e si prende il complemento a
uno del risultato. Il valore è verificato sul disassemblato del gioco~\cite{pokered} e
documentato nella referenza dei formati di questo progetto.

Conviene srotolare la definizione, perché contiene due concetti che il capitolo
precedente ha introdotto e che qui si usano insieme.

Sommare byte in un accumulatore da un byte significa che la somma \emph{tronca}: appena
supera 255 riparte da zero, perché non c'è spazio per il nono bit. In aritmetica
questo si scrive come somma modulo 256, e non è un effetto collaterale da correggere,
è il comportamento su cui l'algoritmo conta.

Il \term{complemento a uno} di un byte è il byte con tutti i bit invertiti, e per un
valore $x$ su otto bit equivale a $255 - x$. Applicarlo alla fine trasforma la somma in
un valore che, sommato ai dati, dà un risultato noto: è una scelta che rende il
controllo verificabile in due modi equivalenti, e la si ritrova in molti protocolli.

Messo in formula, se $b_1 \dots b_n$ sono i byte dell'area protetta:

\[
  \mathrm{checksum} = 255 - \left( \sum_{i=1}^{n} b_i \bmod 256 \right)
\]

Sul gioco, l'area protetta va da \hx{2598} a \hx{3522} e il checksum si trova subito
dopo, a \hx{3523}. La conseguenza di un checksum che non torna è brutale e vale
citarla, perché spiega perché questo progetto tratti la scrittura su cartuccia come
un'operazione irreversibile: il gioco dichiara i dati distrutti e riparte da zero.

Ne segue una regola che vincola qualunque programma voglia modificare un salvataggio di
generazione 1, e che il progetto applica nel proprio codice: dopo ogni modifica il
checksum va ricalcolato, e questo vale sia per un programma che lavora su un file
scaricato dalla cartuccia sia per una routine che scrive direttamente nella memoria
della cartuccia dall'interno del gioco.

\subsection{Quanto è forte, in pratica}

Un checksum a 8 bit costruito come somma ha un difetto che vale conoscere, perché
mostra il limite della categoria. La somma è \emph{commutativa}: non cambia se si
scambiano di posto due byte. Se una corruzione scambia fra loro due byte dell'area
protetta, la somma resta identica e il controllo non se ne accorge. Lo stesso vale per
una corruzione che sottrae a un byte quello che aggiunge a un altro.

È un caso raro fra i guasti fisici, che tendono a azzerare o a saturare interi
settori, ma non è raro affatto fra le manipolazioni deliberate, ed è una delle
ragioni per cui un checksum non è una firma: dice se un dato è probabilmente intatto,
non se è autentico.

\section{Il passo successivo: le parole a sedici bit}

La generazione 3, cioè Rubino, Zaffiro, Smeraldo, Rosso Fuoco e Verde Foglia, usa due
checksum a livelli diversi, e quello che protegge la struttura del singolo Pokemon
merita attenzione per due ragioni: è costruito diversamente, e su di esso le fonti
secondarie sbagliano.

La struttura di un Pokemon in generazione 3 contiene 48 byte cifrati, e il loro
checksum è a 16 bit. La differenza rispetto alla generazione 1 non è solo la
lunghezza: cambia l'unità su cui si somma. Il sorgente del gioco~\cite{pokeemerald}
espone quei 48 byte anche come un vettore di sei parole a 16 bit per ciascuna delle
quattro sottostrutture, e la routine somma quelle parole, non i byte, in un
accumulatore a 16 bit.

\begin{lstlisting}[style=c,caption={La routine di checksum, da include/pokemon.h e src/pokemon.c di pokeemerald~\cite{pokeemerald}. Il tipo dell'accumulatore è la parte che conta.}]
static u16 CalculateBoxMonChecksum(struct BoxPokemon *boxMon)
{
    u16 checksum = 0;
    /* quattro sottostrutture, sei parole da 16 bit ciascuna */
    for (i = 0; i < (s32)ARRAY_COUNT(substruct0->raw); i++)
        checksum += substruct0->raw[i];
    /* idem per substruct1, substruct2, substruct3 */
    return checksum;
}
\end{lstlisting}

Qui c'è l'errore da segnalare, ed è del tipo che distrugge i dati. Una pagina
enciclopedica molto usata~\cite{bulbapedia} descrive questo checksum come una somma
byte per byte. Le due letture danno risultati diversi, e la differenza si vede con un
esempio minimo: se i primi due byte del blocco valgono \hx{FF} e \hx{FF} e tutti gli
altri sono nulli, la somma per parole dà \hx{FFFF}, cioè 65\,535, mentre la somma per
byte dà 510. Chi implementa la seconda scrive un checksum sbagliato, e in generazione 3
un checksum sbagliato non produce un Pokemon strano: lo trasforma in un Uovo Difettoso,
cioè lo distrugge in modo visibile e definitivo.

È il caso che meglio giustifica la gerarchia delle fonti adottata da questo progetto,
per cui il disassemblato del gioco sta sopra qualunque wiki: non perché le wiki siano
inaffidabili in generale, ma perché su un dettaglio come il tipo di un accumulatore la
prosa è ambigua e il codice non lo è.

\subsection{Una proprietà comoda, che sorprende chi la scopre}

C'è un dettaglio di questa costruzione che vale scrivere perché chi lo incontra da
solo rischia di credere di aver trovato un errore. La routine somma le quattro
sottostrutture nel loro ordine logico, ma in memoria quelle quattro sottostrutture sono
permutate in uno di ventiquattro ordini possibili, come il capitolo sulla cifratura
descrive. Poiché l'addizione è commutativa, la somma delle ventiquattro parole non
dipende dall'ordine in cui si prendono i blocchi: ne segue che il checksum si può
verificare senza sapere quale permutazione sia in uso.

È una comodità reale, sfruttata nel codice di questo progetto, ma va dichiarata:
chi legge una verifica del checksum che ignora la permutazione potrebbe pensare che sia
una dimenticanza, e non lo è.

\section{Il terzo livello: proteggere i settori del salvataggio}

Oltre al checksum della singola struttura, la generazione 3 protegge le sezioni del
salvataggio, e l'algoritmo è un terzo tipo ancora. Il salvataggio è diviso in sezioni
da 4096 byte, ciascuna con un piede che contiene un identificativo, un checksum, una
firma costante di valore \hx{08012025} e un indice di salvataggio.

Il checksum di sezione si calcola accumulando le parole da 32 bit della sezione in una
variabile a 32 bit, e poi \emph{ripiegando} il risultato: si sommano i 16 bit alti ai 16
bit bassi, e ciò che resta è il checksum a 16 bit. Il ripiegamento serve a far
contribuire al risultato finale anche i bit alti, che altrimenti andrebbero perduti nel
troncamento, e la stessa tecnica si ritrova nel checksum di Internet.

La firma costante merita una nota, perché è una difesa di natura diversa dal checksum
e i due si confondono facilmente. Il checksum verifica che i dati non siano cambiati; la
firma verifica che quello che si sta leggendo sia effettivamente una sezione di
salvataggio e non un'area di memoria qualunque. Una firma sbagliata invalida la sezione
e con essa l'intero slot, e questo comportamento è documentato nella referenza dei
formati del progetto e verificato sul sorgente~\cite{pokeemerald}.

\subsection{Due copie, e come si decide quale vale}

L'ultima difesa non è un checksum ma un'architettura, e la si capisce partendo dal
guasto che vuole evitare. Se il gioco riscrivesse ogni volta la stessa area, una
mancanza di alimentazione a metà della scrittura lascerebbe quell'area a metà fra il
vecchio e il nuovo, cioè inutilizzabile. La soluzione è non riscrivere mai l'area che
contiene l'ultimo salvataggio valido: si scrive nell'altra.

Il salvataggio di generazione 3 ha quindi due slot da 57\,344 byte, e ciascuna sezione
porta nel piede un indice progressivo. Al caricamento il gioco confronta gli indici e
prende il più recente. Il progetto ha verificato un dettaglio di questo confronto che
conta per chi scrive uno strumento di lettura: lo slot A vince solo se il suo indice è
strettamente maggiore, quindi a parità di indice vince B.

La proprietà che questa architettura garantisce è che in ogni istante esiste almeno
uno slot integro, qualunque cosa accada durante una scrittura. È lo stesso principio
delle scritture atomiche per scambio di puntatore che si usa nei filesystem, ottenuto
qui con due aree fisse e un contatore.

\section{Che cosa ne segue per il progetto}

Da questo capitolo discendono tre vincoli che le altre parti del documento danno per
acquisiti.

Il primo è che qualunque modifica a un salvataggio richiede il ricalcolo di tutti i
checksum che coprono l'area modificata, e che i livelli sono annidati: modificare un
Pokemon in generazione 3 invalida il checksum della sua struttura e quello della sezione
che la contiene.

Il secondo è che un checksum sbagliato non è un avviso ma una distruzione, e cambia
quindi la natura dell'operazione di scrittura. È la ragione tecnica per cui la regola
sull'hardware di questo progetto impone il backup in doppia copia e la rilettura
verificata prima e dopo ogni scrittura su cartuccia, e per cui quella regola non ammette
eccezioni per fretta.

Il terzo è un criterio di lettura delle fonti che il resto del documento applica senza
più ripeterlo: dove una fonte secondaria e un disassemblato divergono su un dettaglio
implementativo, ha ragione il disassemblato. Il caso del checksum a parole contro il
checksum a byte è il primo esempio, e non sarà l'ultimo.

\begin{daverificare}
Il valore della costante di firma e la posizione dei campi del piede di sezione sono
verificati sul sorgente, ma il comportamento del gioco davanti a una firma corretta e un
checksum sbagliato, cioè se invalidi la singola sezione o l'intero slot, è riportato
dalla referenza e non è stato riletto sul sorgente in questa forma. Va confermato prima
di appoggiarvi un recupero.
\end{daverificare}
